\chapter{Énigmes de logique}\label{ch:iv1:logic-riddles}

\begin{ivtrains}
Trouver l'idée unique qui rend une énigme courte. Une pesée donne trois issues, un bit encode un oui ou un non, un coloriage ou une parité est un invariant qu'aucun coup ne peut modifier, et ce que les autres savent est une information. En entretien, énoncer d'abord l'idée, puis la vérifier sur un petit cas.
\end{ivtrains}

\begin{question}[$\star$ \iqroles{QR, QT, Trader, Analyste, QD} \iqfirm{toutes} \iqlevel{junior}]\label{iq:iv1:logic-riddles:1}
Vous disposez de 8 balles d'apparence identique, dont une est légèrement plus lourde, et d'une balance à deux plateaux. Combien de pesées faut-il pour trouver à coup sûr la balle lourde, et comment~? Combien de balles pourriez-vous traiter avec 4 pesées~?
\end{question}

\begin{question}[$\star$ \iqroles{QR, QT, Trader, Analyste} \iqfirm{toutes} \iqlevel{junior}]\label{iq:iv1:logic-riddles:2}
Vous disposez de deux cordes et d'un briquet. Chaque corde met exactement 60 minutes à brûler d'un bout à l'autre, mais elle brûle de façon irrégulière, si bien que la moitié d'une corde ne met pas forcément 30 minutes. Comment mesurer exactement 45 minutes~?
\end{question}

\begin{question}[$\star$ \iqroles{QR, QT, Trader, Analyste} \iqfirm{toutes} \iqlevel{junior}]\label{iq:iv1:logic-riddles:3}
Vous disposez d'un sablier de 4 minutes et d'un sablier de 7 minutes, et vous pouvez retourner l'un ou l'autre au moment où l'un d'eux se vide. Comment mesurer exactement 9 minutes, en les démarrant tous deux à l'instant 0~?
\end{question}

\begin{question}[$\star$ \iqroles{QR, QT, Trader, Analyste, QD} \iqfirm{toutes} \iqlevel{junior}]\label{iq:iv1:logic-riddles:4}
Il y a 10 sacs de pièces. Dans neuf d'entre eux, chaque pièce pèse 10 grammes, et dans l'un d'eux, chaque pièce pèse 11 grammes. Vous disposez d'une balance exacte qui donne un nombre, et vous pouvez l'utiliser une seule fois. Comment trouver le sac~?
\end{question}

\begin{question}[$\star$ \iqroles{QR, QT, Analyste} \iqfirm{toutes} \iqlevel{junior}]\label{iq:iv1:logic-riddles:5}
Vous devez prendre exactement un comprimé de type A et un comprimé de type B chaque jour. Les comprimés ont la même apparence, et ils sont chers. Aujourd'hui, vous avez fait tomber un A dans votre main, puis, par accident, deux B, et vous ne pouvez pas distinguer les trois. Comment obtenir la dose correcte du jour sans gaspiller un comprimé~?
\end{question}

\begin{question}[$\star$ \iqroles{QR, QT, Analyste, QD} \iqfirm{toutes} \iqlevel{junior}]\label{iq:iv1:logic-riddles:6}
Une chaîne en or compte 7 maillons alignés. Vous devez payer un ouvrier à raison d'un maillon d'or à la fin de chacun des 7 jours, et vous pouvez récupérer des morceaux précédemment remis en monnaie. Couper un maillon le libère sous forme de pièce unique et découpe la chaîne de part et d'autre. Quel est le nombre minimal de maillons à couper, et lesquels~? Avec 2 coupures, quelle est la chaîne la plus longue que vous pourriez payer de cette façon, un maillon par jour~?
\end{question}

\begin{question}[$\star$ \iqroles{QR, QT, Trader, Analyste} \iqfirm{toutes} \iqlevel{junior}]\label{iq:iv1:logic-riddles:7}
Deux cases d'angle opposées sont retirées d'un échiquier de 8 sur 8. Les 62 cases restantes peuvent-elles être recouvertes par 31 dominos, chacun couvrant deux cases voisines~? Et si vous retirez plutôt deux cases de couleurs opposées, n'importe où sur le plateau~?
\end{question}

\begin{question}[$\star$ \iqroles{QR, QT, Trader, Analyste} \iqfirm{toutes} \iqlevel{junior}]\label{iq:iv1:logic-riddles:8}
Deux personnes reçoivent chacune au hasard un chapeau noir ou blanc. Chacune ne voit que le chapeau de l'autre, et elles annoncent leur propre couleur au même moment, sans aucun signal. Elles peuvent convenir d'un plan au préalable. Elles gagnent si au moins une des deux réponses est juste. Trouver un plan qui gagne toujours.
\end{question}

\begin{question}[$\star$ \iqroles{QR, QT, Analyste} \iqfirm{toutes} \iqlevel{junior}]\label{iq:iv1:logic-riddles:9}
Trois interrupteurs placés à l'extérieur d'une pièce fermée commandent trois ampoules à filament ordinaires à l'intérieur. Vous pouvez régler les interrupteurs comme vous le souhaitez, puis entrer dans la pièce une seule fois, après quoi vous ne pouvez plus les toucher. Comment savoir quel interrupteur commande quelle ampoule~?
\end{question}

\begin{question}[$\star$ \iqroles{QR, QT, Analyste, Trader} \iqfirm{toutes} \iqlevel{junior}]\label{iq:iv1:logic-riddles:10}
Trois boîtes sont étiquetées «~pommes~», «~oranges~» et «~mixte~», et on vous dit que chaque étiquette est fausse. Vous pouvez prendre un fruit dans une boîte de votre choix sans regarder à l'intérieur. Quelle boîte choisissez-vous, et comment réétiquetez-vous alors les trois~?
\end{question}

\begin{question}[$\star$ \iqroles{QR, QT, Trader, Analyste} \iqfirm{toutes} \iqlevel{junior}]\label{iq:iv1:logic-riddles:11}
Quatre personnes doivent traverser un pont de nuit avec une seule lampe torche. Deux personnes au plus peuvent traverser à la fois, et toute traversée nécessite la lampe, qu'il faut ramener à pied. Elles mettent 1, 2, 5 et 10 minutes pour traverser, et un duo avance à l'allure du plus lent. Quel est le temps total minimal, et quel est le plan~?
\end{question}

\begin{question}[$\star$ \iqroles{QR, QT, Analyste} \iqfirm{toutes} \iqlevel{junior}]\label{iq:iv1:logic-riddles:12}
Vous disposez d'un bidon de 3 litres, d'un bidon de 5 litres et d'une quantité d'eau illimitée, sans autre repère. Comment mesurer exactement 4 litres~? Pourriez-vous mesurer 3 litres avec un bidon de 6 litres et un bidon de 10 litres~? Énoncer la règle générale.
\end{question}

\begin{question}[$\star$ \iqroles{QR, QT, QD} \iqfirm{toutes} \iqlevel{junior}]\label{iq:iv1:logic-riddles:13}
Sur une grille de 8 sur 8, une case est noire et toutes les autres sont blanches. Un coup inverse toutes les cases d'une ligne ou toutes les cases d'une colonne. Peut-on atteindre une grille entièrement blanche~? Quel est l'invariant~?
\end{question}

\begin{question}[$\star\star$ \iqroles{QR, QT, Trader, Analyste, QD} \iqfirm{toutes} \iqlevel{junior}]\label{iq:iv1:logic-riddles:14}
Cent personnes se tiennent en file, chacune avec un chapeau noir ou blanc tiré au hasard, et chacune ne voit que les chapeaux situés devant elle. Depuis le fond de la file, chacune à son tour annonce «~noir~» ou «~blanc~», et tout le monde entend. Elles peuvent convenir d'un plan au préalable. Combien peuvent être certaines d'avoir raison, et comment~?
\end{question}

\begin{question}[$\star\star$ \iqroles{QR, QT, QD, Dev} \iqfirm{toutes} \iqlevel{junior}]\label{iq:iv1:logic-riddles:15}
L'une de 1000 bouteilles de vin a été empoisonnée. Un dégustateur présente l'effet exactement 24 heures après une gorgée, quelle que soit sa taille, et vous devez connaître la bouteille dans 24 heures. Combien de dégustateurs faut-il au minimum~? Combien en faudrait-il avec 48 heures, et en pouvant utiliser un second tour avec ceux qui ne sont pas encore touchés~?
\end{question}

\begin{question}[$\star\star$ \iqroles{QR, QT, Trader} \iqfirm{toutes} \iqlevel{junior}]\label{iq:iv1:logic-riddles:16}
Une table porte 100 pièces, dont on vous dit que exactement 10 montrent face. Vous avez les yeux bandés et vous ne pouvez pas sentir quel côté est dessus. Séparer les pièces en deux groupes, en retournant les pièces de votre choix, de sorte que les deux groupes comptent le même nombre de faces.
\end{question}

\begin{question}[$\star\star$ \iqroles{QR, QT, Analyste} \iqfirm{toutes} \iqlevel{junior}]\label{iq:iv1:logic-riddles:17}
Un cavalier peut-il partir d'un coin d'un échiquier de 8 sur 8, visiter chaque case exactement une fois, et finir dans le coin opposé~? Donner l'argument. Pourrait-il finir sur une autre case, en principe~?
\end{question}

\begin{question}[$\star\star$ \iqroles{QR, QT, Analyste} \iqfirm{toutes} \iqlevel{junior}]\label{iq:iv1:logic-riddles:18}
Une case est retirée d'un échiquier de $2^n$ sur $2^n$, n'importe laquelle. Montrer que le reste peut être recouvert par des pièces en forme de L de trois cases. Combien de pièces utilise-t-on sur un échiquier de 8 sur 8~?
\end{question}

\begin{question}[$\star\star$ \iqroles{QR, QT, Trader} \iqfirm{toutes} \iqlevel{junior}]\label{iq:iv1:logic-riddles:19}
Une île compte 13 caméléons rouges, 15 verts et 17 bleus. Quand deux caméléons de couleurs différentes se rencontrent, tous deux changent pour la troisième couleur. Tous les caméléons peuvent-ils un jour être de la même couleur~? Et si les effectifs étaient 13, 16 et 16~?
\end{question}

\begin{question}[$\star\star$ \iqroles{QR, QT, QD} \iqfirm{toutes} \iqlevel{junior}]\label{iq:iv1:logic-riddles:20}
Les nombres de 1 à 10 sont écrits au tableau. Un coup efface deux nombres quelconques $a$ et $b$ et écrit $|a - b|$. Après neuf coups, il ne reste qu'un nombre. Peut-il valoir 0~? Par ailleurs, si chaque coup écrit plutôt $a + b + ab$, quel est le dernier nombre, et l'ordre compte-t-il~?
\end{question}

\begin{question}[$\star\star$ \iqroles{QR, QT, QD, Dev} \iqfirm{toutes} \iqlevel{junior}]\label{iq:iv1:logic-riddles:21}
Vous disposez de 100 nombres distincts. Combien de comparaisons faut-il pour trouver à la fois le plus grand et le plus petit, en procédant astucieusement, plutôt que les $198$ de deux parcours séparés~? Donner la méthode, et la formule pour $n$ nombres.
\end{question}

\begin{question}[$\star\star$ \iqroles{QR, QT, Trader, Analyste} \iqfirm{toutes} \iqlevel{junior}]\label{iq:iv1:logic-riddles:22}
Un recenseur demande à une femme l'âge de ses trois filles, en nombres entiers. Elle répond que le produit des âges vaut 36, et que la somme est le numéro de la maison voisine. Le recenseur regarde ce numéro et ne peut toujours pas conclure. Elle ajoute que l'aînée joue du piano. Quels sont les âges, et quel doit être le numéro de la maison~?
\end{question}

\begin{question}[$\star\star$ \iqroles{QR, QT, Trader, Analyste} \iqfirm{toutes} \iqlevel{junior}]\label{iq:iv1:logic-riddles:23}
Trois logiciens parfaits se tiennent en file, chacun voyant les chapeaux situés devant lui. Les chapeaux ont été tirés d'une boîte contenant 3 rouges et 2 blancs. Celui du fond dit~: «~Je ne connais pas ma couleur.~» Celui du milieu dit~: «~Moi non plus.~» Celui de devant, qui ne voit personne, dit~: «~Je connais la mienne.~» Quelle est la couleur, et comment le sait-il~?
\end{question}

\begin{question}[$\star\star$ \iqroles{QR, QT, QD} \iqfirm{toutes} \iqlevel{junior}]\label{iq:iv1:logic-riddles:24}
Il y a 10 piles de pièces. Chaque pile est soit entièrement authentique, à 10 grammes la pièce, soit entièrement fausse, à 11 grammes la pièce, et un nombre quelconque de piles peut être faux. Vous pouvez utiliser une balance exacte une seule fois. Comment trouver toutes les piles fausses~? Combien de pièces faut-il en tout~?
\end{question}

\begin{question}[$\star\star$ \iqroles{QR, QT, Analyste} \iqfirm{toutes} \iqlevel{junior}]\label{iq:iv1:logic-riddles:25}
Vous vous tenez à une bifurcation, l'une des routes menant au salut, l'autre non. Deux gardes connaissent la route. L'un dit toujours la vérité et l'autre ment toujours, et vous ne savez pas lequel est lequel. Avec une seule question à réponse oui ou non posée à un garde, comment trouver la route sûre~? Montrer que votre question fonctionne dans les quatre cas.
\end{question}

\begin{question}[$\star\star$ \iqroles{QR, QT, Analyste} \iqfirm{toutes} \iqlevel{junior}]\label{iq:iv1:logic-riddles:26}
Un carnet compte 100 énoncés numérotés. L'énoncé $k$ dit~: «~Exactement $k$ des énoncés de ce carnet sont faux.~» Quels énoncés sont vrais~? Votre réponse changerait-elle si l'énoncé disait «~au moins $k$~»~?
\end{question}

\begin{question}[$\star\star\star$ \iqroles{QR, QT, QD, Trader} \iqfirm{toutes} \iqlevel{junior, confirmé}]\label{iq:iv1:logic-riddles:27}
Vous disposez de 12 boules, dont exactement une est anormale~: plus lourde ou plus légère, sans que vous sachiez laquelle. Avec une balance, trouvez la boule anormale et dites si elle est plus lourde ou plus légère, en 3 pesées. Pourquoi 13 boules sont-elles impossibles sous cette forme, alors que $3^3 = 27$ issues semblent suffire~? Donnez un schéma de pesées fixé à l'avance.
\end{question}

\begin{question}[$\star\star\star$ \iqroles{QR, QT, QD} \iqfirm{toutes} \iqlevel{junior, confirmé}]\label{iq:iv1:logic-riddles:28}
Pouvez-vous recouvrir exactement un plateau de 10 sur 10 avec des briques droites de 1 sur 4~? Qu'en est-il d'un plateau de 4 sur 10, et quelle est la règle générale pour un plateau $m$ sur $n$ avec des briques de 1 sur 4~? Donnez le coloriage qui tranche.
\end{question}

\begin{question}[$\star\star\star$ \iqroles{QR, QT, QD, Dev} \iqfirm{toutes} \iqlevel{junior, confirmé}]\label{iq:iv1:logic-riddles:29}
Combien de comparaisons faut-il pour trouver le deuxième plus grand de 64 nombres distincts, par rapport à deux parcours séparés~? Donnez la méthode et la formule pour $n$ nombres. Pourquoi tous les nombres sauf deux doivent-ils perdre une comparaison face à un autre que le plus grand~?
\end{question}

\begin{question}[$\star\star\star$ \iqroles{QR, QT, Trader} \iqfirm{hedge fund, teneur de marché} \iqlevel{junior, confirmé}]\label{iq:iv1:logic-riddles:30}
Un groupe de $n$ personnes reçoit chacune un chapeau, dont la couleur est choisie parmi $n$ couleurs, les répétitions étant permises. Chacune voit tous les chapeaux sauf le sien, et elles devinent toutes au même instant. Elles peuvent convenir d'un plan au préalable. Trouvez un plan qui garantit qu'exactement une personne a raison, quels que soient les chapeaux.
\end{question}

\begin{question}[$\star\star\star$ \iqroles{QR, QT, Trader} \iqfirm{hedge fund} \iqlevel{junior, confirmé}]\label{iq:iv1:logic-riddles:31}
Trois personnes reçoivent chacune un chapeau noir ou blanc, tiré au hasard. Chacune voit les deux autres chapeaux et, au même instant, elles annoncent leur couleur ou passent. Le groupe gagne si au moins une personne annonce une couleur et si toutes les annonces sont justes. Elles peuvent convenir d'un plan au préalable. Quelle probabilité de gagner peuvent-elles atteindre, et pourquoi cela ne contredit-il pas la règle selon laquelle chaque annonce est juste une fois sur deux~?
\end{question}

\begin{question}[$\star\star\star$ \iqroles{QR, QT, Trader} \iqfirm{hedge fund, teneur de marché} \iqlevel{junior, confirmé}]\label{iq:iv1:logic-riddles:32}
Sur une île vivent 100 personnes aux yeux bleus et 100 aux yeux marron, et personne ne sait que ce sont les deux seules couleurs possibles. Personne ne peut discuter de la couleur des yeux, il n'y a pas de miroirs, et quiconque déduit la couleur de ses propres yeux quitte l'île à minuit ce jour-là. Un visiteur annonce à tous~: «~Au moins l'un d'entre vous a les yeux bleus.~» Que se passe-t-il~? Quelle information nouvelle le visiteur a-t-il donnée, alors que chacun voyait déjà des yeux bleus~?
\end{question}

\begin{question}[$\star\star\star$ \iqroles{QR, QT, DS, QD} \iqfirm{toutes} \iqlevel{junior, confirmé}]\label{iq:iv1:logic-riddles:33}
Cent prisonniers peuvent se concerter une seule fois, puis sont isolés les uns des autres. Chaque jour, un prisonnier choisi au hasard entre dans une pièce munie d'un interrupteur, éteint au départ, et peut le basculer. Tout prisonnier peut déclarer «~nous sommes tous passés par la pièce~». Si cette déclaration est juste, ils sont libérés~; si elle est fausse, ils restent tous enfermés. Trouvez un plan qui fonctionne à coup sûr. Combien de temps faut-il en moyenne~?
\end{question}

\begin{question}[$\star\star\star\star$ \iqroles{QR, QT, QD} \iqfirm{hedge fund, teneur de marché} \iqlevel{confirmé}]\label{iq:iv1:logic-riddles:34}
Un échiquier porte une pièce sur chacune de ses 64 cases, côté pile ou côté face, dans une disposition choisie par un geôlier hostile, qui désigne aussi une case. Le prisonnier A voit tout cela et doit retourner exactement une pièce. Le prisonnier B entre ensuite, ne voit que les pièces et doit nommer la case que le geôlier a choisie. Ils peuvent se concerter au préalable. Trouvez le plan, et dites ce qu'il y a de spécial dans le nombre 64.
\end{question}

\begin{question}[$\star\star\star\star$ \iqroles{QR, QT, Trader} \iqfirm{hedge fund} \iqlevel{confirmé}]\label{iq:iv1:logic-riddles:35}
Un spectateur choisit 5 cartes dans un jeu de cartes standard et les remet à un assistant. L'assistant rend l'une d'elles au spectateur et dispose les 4 autres en ligne, face visible, dans l'ordre de son choix. Un magicien, qui n'a rien vu d'autre, nomme la carte cachée. Comment fonctionne le système~? Faites un argument de comptage pour montrer jusqu'où ce nombre pourrait dépasser 52.
\end{question}
